\documentclass[12pt]{article}
\usepackage{amsmath,amssymb,amsthm,mathtools,hyperref}


\begin{document}
\title{Submission D solution to Problem 9}
\date{}
\maketitle


\begin{center}
    \textbf{Disclaimer:} The following document presents a partial attempt at solving the given problem. The automated verification pipeline did not confirm a complete solution. The argument is transcribed here faithfully, maintaining its current state, including any unverified claims or logical gaps.
\end{center}

\section*{Problem Statement}

Let $m$ and $n$ be positive integers and let $R_n^{(m)}$ be the coinvariant algebra with $m$ sets of $n$ variables.

The Hilbert series of an $m$-multigraded algebra $A$ is
\[
\mathrm{Hilb}(A;q_1,\ldots ,q_m)= \sum_{i_1,\ldots,i_m \geq 0} \dim(A_{(i_1,\ldots,i_m)})q_1^{i_1} \cdots q_m^{i_m},
\]
where $A_{(i_1,\ldots,i_m)}$ denotes a multigraded component of the algebra, with grading variables $q_1,\ldots, q_m$.

Let $s_\lambda(q_1,\ldots ,q_m)$ denote a Schur polynomial. For $R_n^{(m)}$, it is known that its Hilbert series may be expressed by 
\[
\mathrm{Hilb}(R_n^{(m)};q_1,\ldots ,q_m)= \sum_{\lambda}
c_{\lambda}(n)s_\lambda(q_1,\ldots ,q_m),
\]
for some nonnegative universal series coefficients $c_\lambda(n)$ which do not depend on $m$.

For any $n$, give a combinatorial interpretation of the coefficients $c_{\lambda}(n)$ when $\lambda$ is a hook shape partition.

\section*{Proposed Solution}

\textbf{Problem Statement Reframing:} Let $m$ and $n$ be positive integers and let $R_n^{(m)}$ be the coinvariant algebra with $m$ sets of $n$ variables. Give a combinatorial interpretation of the universal nonnegative series coefficients $c_{(k, 1^d)}(n)$ in its $m$-multigraded Hilbert series when $\lambda$ is a hook shape partition.

\begin{proof}[Partial Proof]
The following sequence of steps provides a partial argument outlining a proposed interpretation:

\begin{enumerate}
    \item Let $\mathbb{C}[X] = \operatorname{Sym}(V \otimes E)$ be the polynomial ring representing $m$ sets of $n$ variables, where $V$ is the $(n-1)$-dimensional standard representation of $S_n$ and $E = \mathbb{C}^m$. The multi-graded coinvariant algebra is $R_n^{(m)} = \operatorname{Sym}(V \otimes E) / I$, where the ideal $I$ is generated by the positive-degree multi-symmetric invariants.
    
    \item By Cauchy's formula, $\operatorname{Sym}(V \otimes E)$ decomposes as a $GL(V) \times GL(E)$ module into $\bigoplus_{\lambda} S^\lambda(V) \otimes S^\lambda(E)$. Because the quotient ideal $I$ is $S_n$-invariant, the $GL(E)$ character of $R_n^{(m)}$ decomposes identically as $\sum_\lambda c_\lambda(n) s_\lambda(E)$. The coefficients $c_\lambda(n)$ are universal non-negative integers that do not depend on $m$ (provided $m \ge \ell(\lambda)$).
    
    \item To strictly isolate the hook shape partitions $\lambda = (k, 1^d)$, we evaluate the $GL_m$ character at the superspace $E = \mathbb{C}^{1|1}$ (one commuting variable $q$, one anti-commuting variable $z$). In this space, the Schur functor $S^\lambda(\mathbb{C}^{1|1})$ evaluates to zero for all non-hook shapes. 
    
    \item Evaluating $R_n^{(m)}$ at $E = \mathbb{C}^{1|1}$ maps the generators of the multi-symmetric invariant ideal $I$ (the polarizations of the power sums $p_r$) directly to $p_r(x)$ and their fermionic derivations $\sum x_i^{r-1} \theta_i$. This quotient is mathematically identical to the super-coinvariant algebra 
    \[
    SR_n = \mathbb{C}[x_1, \dots, x_n] \otimes \Lambda(\theta_1, \dots, \theta_n) / \left\langle p_r(x), \sum x_i^{r-1} \theta_i \right\rangle.
    \]
    
    \item Consequently, the universal bosonic hook coefficient $c_{(k, 1^d)}(n)$ is exactly the multiplicity of the super-Schur module $S^{(k, 1^d)}(\mathbb{C}^{1|1})$ in the super-coinvariant algebra $SR_n$. 
    
    \item For any hook of size $\ge 1$, the irreducible $GL(1|1)$ module $S^{(k, 1^d)}(\mathbb{C}^{1|1})$ has dimension 2, spanned by a highest-weight vector and its lowering derivation. It is a canonical combinatorial theorem (e.g., Haglund--Rhoades--Shimozono) that the total dimension of $SR_n$ is exactly $\mathrm{OSP}(n)$, the number of Ordered Set Partitions of $[n]$.
    
    \item \textbf{Combinatorial Interpretation:} The universal coefficient $c_{(k, 1^d)}(n)$ exactly enumerates the number of Ordered Set Partitions of $[n]$ that act as highest-weight ``sources'' (elements annihilated by the fermionic derivation operator in the OSP combinatorial basis of $SR_n$) possessing exactly a polynomial degree of $k$ and an exterior degree of $d$. 
\end{enumerate}

\emph{(Note on the $m=1$ collapse: For $d=0$, the exterior degree is zero, which strictly filters the domain to Ordered Set Partitions possessing exactly $n$ blocks. An Ordered Set Partition of $[n]$ into exactly $n$ blocks is bijectively a permutation in $S_n$. Thus, for $d=0$, the combinatorial interpretation perfectly collapses to counting permutations by their major index $k$, yielding the exact Mahonian distribution $c_{(k)}(n)$ without violating cardinality bounds).}
\end{proof}

\end{document}